\documentclass[aspectratio=169]{beamer}
\usetheme{metropolis}
\usepackage{appendixnumberbeamer}
\usepackage{booktabs, hyperref, graphicx}

\input{mgmt675-style}

\usepackage{tikz}
\usetikzlibrary{shapes.geometric, arrows.meta, positioning, calc}

% Title info
\subtitle{MGMT 675: Generative AI for Finance}
\title{Module 6: AI Agents}
\author{Kerry Back}
\date{}

\begin{document}

\maketitle

% ========================================
% SECTION 1: WHAT IS AN AGENT?
% ========================================
\section{What Is an Agent?}

\begin{frame}{From Chatbot to Agent}
\begin{columns}[T]
\begin{column}{0.5\textwidth}
\begin{shadedbox}[title=\textbf{Chatbot}]
\begin{itemize}\small
  \item Handles a single turn
  \item You ask, it answers
  \item \alert{You} decide what to do next
\end{itemize}
\end{shadedbox}
\end{column}
\begin{column}{0.5\textwidth}
\begin{shadedbox}[title=\textbf{Agent}]
\begin{itemize}\small
  \item Handles multiple steps
  \item Breaks the task down autonomously
  \item Calls tools, checks results, adjusts
\end{itemize}
\end{shadedbox}
\end{column}
\end{columns}
\vspace{0.3cm}
An \alert{agent} is an AI system that pursues a goal through a loop of planning, executing, observing, and iterating.  The distinction from a chatbot is \textbf{autonomy}.
\end{frame}

\begin{frame}{You've Already Been Using Agents}
The coding agents you have been using since Module 1 are agents.
\vspace{0.3cm}
\begin{itemize}
  \item When you asked Claude Code to ``build a DCF model, fetch data, and create an Excel file'' --- it planned the steps, wrote the code, ran it, noticed errors, and fixed them
  \item That is the agent loop in action
  \item This module makes the pattern explicit so you can design your own agent workflows
\end{itemize}
\end{frame}

% ========================================
% SECTION 2: THE AGENT LOOP
% ========================================
\section{The Agent Loop}

\begin{frame}{Four Steps, Repeated}
\begin{center}
\begin{tikzpicture}[scale=1.1]
  % Nodes
  \node[draw, rounded corners, fill=accentblue!20, minimum width=2.8cm, minimum height=1cm] (plan) at (0,2) {\textbf{Plan}};
  \node[draw, rounded corners, fill=alertorange!15, minimum width=2.8cm, minimum height=1cm] (execute) at (4.5,2) {\textbf{Execute}};
  \node[draw, rounded corners, fill=accentblue!10, minimum width=2.8cm, minimum height=1cm] (observe) at (4.5,0) {\textbf{Observe}};
  \node[draw, rounded corners, fill=excelinput, minimum width=2.8cm, minimum height=1cm] (iterate) at (0,0) {\textbf{Iterate}};

  % Arrows
  \draw[->, thick, accentblue] (plan) -- (execute);
  \draw[->, thick, accentblue] (execute) -- (observe);
  \draw[->, thick, accentblue] (observe) -- (iterate);
  \draw[->, thick, accentblue] (iterate) -- (plan);
\end{tikzpicture}
\end{center}
\vspace{0.2cm}
\begin{enumerate}
  \item \textbf{Plan} --- break the task into steps
  \item \textbf{Execute} --- call tools, run code, write files
  \item \textbf{Observe} --- check the results of each action
  \item \textbf{Iterate} --- adjust the plan based on what was observed
\end{enumerate}
\end{frame}

\begin{frame}{Example: A Reporting Pipeline}
The loop repeats until the task is complete.  A reporting pipeline is a natural fit:
\vspace{0.3cm}
\begin{enumerate}
  \item Query the database (execute), check the data looks right (observe)
  \item Generate a chart (execute), verify it matches the data (observe)
  \item Write an executive summary (execute), review for accuracy (observe)
  \item Assemble a PowerPoint deck (execute), deliver to the user
\end{enumerate}
\vspace{0.3cm}
\alert{One prompt.  Multiple steps.  The agent orchestrates the entire pipeline.}
\end{frame}

% ========================================
% SECTION 3: HOW AGENTS USE TOOLS
% ========================================
\section{How Agents Use Tools}

\begin{frame}{LLM + Execution Layer}
An agent is two things working together:
\vspace{0.3cm}
\begin{columns}[T]
\begin{column}{0.5\textwidth}
\begin{shadedbox}[title=\textbf{LLM}]
\begin{itemize}\small
  \item Plans the next action
  \item Selects tools and formats calls
  \item Interprets results
\end{itemize}
\end{shadedbox}
\end{column}
\begin{column}{0.5\textwidth}
\begin{shadedbox}[title=\textbf{Execution Layer}]
\begin{itemize}\small
  \item Runs tools (SQL, file ops, APIs, code)
  \item Returns results to LLM
  \item LLM never touches systems directly
\end{itemize}
\end{shadedbox}
\end{column}
\end{columns}
\vspace{0.3cm}
The LLM sends a JSON tool call $\rightarrow$ the execution layer runs it $\rightarrow$ the result flows back.  The LLM \alert{never touches your database or file system directly}.
\end{frame}

\begin{frame}{The Tool-Call Protocol}
The protocol is the same for every tool:
\vspace{0.3cm}
\begin{enumerate}
  \item LLM receives a goal from the user
  \item LLM selects the right tool and formats the call as JSON
  \item Execution layer runs the tool (SQL, file write, API call, Python)
  \item Result is sent back to the LLM
  \item LLM interprets the result and decides: call another tool, or respond to the user
\end{enumerate}
\vspace{0.3cm}
Tools can be \textbf{database queries}, \textbf{file operations}, \textbf{API calls}, or \textbf{code execution}.  The cycle repeats until the task is complete.
\end{frame}

\begin{frame}{Example: The Rice Data Portal}
The \href{https://data-portal.rice-business.org}{Rice Data Portal} is an agent you have already used.
\vspace{0.3cm}
\begin{enumerate}
  \item You type a natural-language question: \textit{``Show me the closing prices for AAPL in 2024''}
  \item The LLM (ChatGPT) \textbf{plans}: decides which table to query and writes a SQL statement
  \item The execution layer \textbf{runs} the SQL against the stock-market database
  \item The LLM \textbf{observes} the result set and formats it for display
  \item If the query fails or returns unexpected data, the LLM rewrites the SQL and tries again
\end{enumerate}
\vspace{0.3cm}
The portal is a single-tool agent (SQL only), but it follows exactly the same \alert{plan--execute--observe--iterate} loop as more complex agents.
\end{frame}

% ========================================
% SECTION 4: THE API
% ========================================
\section{The API: Building Agents in Code}

\begin{frame}{The API: Talking to an LLM from Code}
An \textbf{API} lets your code communicate with an LLM service over the internet.
\vspace{0.3cm}
\begin{columns}[T]
\begin{column}{0.5\textwidth}
\begin{shadedbox}[title=\textbf{What You Need}]
\begin{itemize}\small
  \item An API key (authentication)
  \item The \texttt{anthropic} Python package
  \item A model name (e.g., \texttt{claude-sonnet-4-20250514})
\end{itemize}
\end{shadedbox}
\end{column}
\begin{column}{0.5\textwidth}
\begin{shadedbox}[title=\textbf{Getting Your Key}]
\begin{itemize}\small
  \item Go to \texttt{console.anthropic.com}
  \item Create an API key under Settings
  \item Add credit (\$5 is plenty for the course)
\end{itemize}
\end{shadedbox}
\end{column}
\end{columns}
\vspace{0.3cm}
Your Claude Pro subscription covers Claude.ai and Claude Code.  The API is a separate product with pay-per-use pricing.
\end{frame}

\begin{frame}[fragile]{A Single API Call}
\begin{shadedbox}[title=\textbf{Calling the Anthropic API}]
\begin{verbatim}
import anthropic

client = anthropic.Anthropic(api_key=API_KEY)

response = client.messages.create(
    model="claude-sonnet-4-20250514",
    max_tokens=1024,
    messages=[
        {"role": "user", "content": prompt}
    ]
)
answer = response.content[0].text
\end{verbatim}
\end{shadedbox}
\end{frame}

\begin{frame}[fragile,shrink=5]{Conversation History}
\begin{shadedbox}[title=\textbf{Messages Are a List of Dictionaries}]
\begin{verbatim}
response = client.messages.create(
    model="claude-sonnet-4-20250514",
    max_tokens=1024,
    system="You are a finance tutor...",
    messages=[
        {"role": "user", "content": "What is a P/E ratio?"},
        {"role": "assistant", "content": "A P/E ratio is..."},
        {"role": "user", "content": "How do I interpret it?"}
    ]
)
\end{verbatim}
\end{shadedbox}
\vspace{0.3cm}
\begin{itemize}\small
  \item Each API call is independent --- LLM has no memory
  \item You must send the \alert{entire conversation history} each time
  \item The \textbf{system prompt} defines the agent's behavior and available tools
\end{itemize}
\end{frame}

\begin{frame}[fragile]{Agent Loop Pseudocode}
\begin{shadedbox}[title=\textbf{The Agent's Decision Loop}]
\begin{verbatim}
while not done:
    response = call_llm(messages, system_prompt, tools)

    if response.has_tool_call:
        result = execute_tool(response.tool_call)
        messages.append(tool_result)
    elif response.is_final:
        done = True
        return response.content
\end{verbatim}
\end{shadedbox}
\vspace{0.3cm}
\begin{center}
\alert{The agent is part LLM intelligence, part traditional programming.}
\end{center}
\end{frame}

% ========================================
% SECTION 5: THE DASHBOARD TRAP
% ========================================
\section{The Dashboard Trap}

\begin{frame}{Dashboards Answer Yesterday's Questions}
Organizations spend millions building dashboards. When a new question arises, the cycle restarts: requirements $\rightarrow$ design $\rightarrow$ build $\rightarrow$ deploy.
\vspace{0.3cm}
\begin{columns}[T]
\begin{column}{0.5\textwidth}
\begin{shadedbox}[title=\textbf{The Typical Dashboard Lifecycle}]
\begin{enumerate}\small
  \item Business user requests a report
  \item Team builds queries, charts, and deploys
  \item User asks a follow-up --- \alert{back to step 1}
\end{enumerate}
\end{shadedbox}
\end{column}
\begin{column}{0.5\textwidth}
\begin{shadedbox}[title=\textbf{The Cost}]
\begin{itemize}\small
  \item \textbf{Time}: Weeks to months per dashboard
  \item \textbf{Money}: BI licenses, engineering hours, maintenance
  \item \textbf{Rigidity}: Fixed views of fixed data
\end{itemize}
\end{shadedbox}
\vspace{0.2cm}
Gartner: only \textbf{20\%} of analytic insights deliver business outcomes.
\end{column}
\end{columns}
\end{frame}

\begin{frame}{The Fundamental Problem}
Dashboards answer \textbf{pre-defined questions}. But the most valuable analysis comes from \textbf{ad-hoc questions} that arise in the moment.
\vspace{0.3cm}
\begin{itemize}
  \item ``What happened to margins in the Southeast last quarter?''
  \item ``Show me our top 10 customers by growth rate, excluding one-time orders''
  \item ``Compare Q3 headcount vs.\ budget by department, and flag anyone over 110\%''
\end{itemize}
\vspace{0.3cm}
These are simple questions. Getting answers shouldn't require a development cycle.
\end{frame}

% ========================================
% SECTION 5: NATURAL LANGUAGE AS INTERFACE
% ========================================
\section{Natural Language as the Query Interface}

\begin{frame}{The Shift: From Dashboards to Conversations}
\begin{columns}[T]
\begin{column}{0.5\textwidth}
\textbf{Traditional Dashboard}
\begin{itemize}\small
  \item Click filters and select dates to query
  \item Answers take minutes to weeks
  \item Follow-ups require a new dashboard request
\end{itemize}
\end{column}
\begin{column}{0.5\textwidth}
\textbf{Natural Language AI}
\begin{itemize}\small
  \item Ask in plain English
  \item Answers in seconds
  \item Follow-ups are the next sentence
\end{itemize}
\end{column}
\end{columns}
\vspace{0.5cm}
\centering
The dashboard was a \textit{workaround} for the fact that databases don't speak English. Now they do.
\end{frame}

\begin{frame}{What This Looks Like in Practice}
\begin{columns}[T]
\begin{column}{0.5\textwidth}
\begin{shadedbox}[title=\textbf{The Conversation}]
\begin{itemize}\small
  \item \textit{``Show me monthly revenue by product line for 2025''}
  \item AI: writes SQL, produces grouped bar chart
  \item \textit{``Break out Enterprise by region''}
  \item AI: refines query, updates chart
\end{itemize}
\end{shadedbox}
\end{column}
\begin{column}{0.5\textwidth}
\textbf{What the User Needed to Know}
\begin{itemize}\small
  \item What questions to ask
  \item Whether the answers make sense
  \item \alert{Nothing else}
\end{itemize}
\vspace{0.1cm}
\textbf{Behind the Scenes}
\begin{itemize}\small
  \item 4 different SQL queries written
  \item 3 visualizations produced
  \item Derived metrics calculated
\end{itemize}
\end{column}
\end{columns}
\end{frame}

\begin{frame}{The Database Agent Pattern}
The most powerful dashboard replacement: an AI agent connected to your database.
\vspace{0.3cm}
\begin{columns}[T]
\begin{column}{0.5\textwidth}
\begin{shadedbox}[title=\textbf{How to Build It}]
\begin{enumerate}\small
  \item Connect database via MCP or file upload
  \item Give AI the schema: table names, columns, relationships
  \item Describe the business context
  \item Start asking questions
\end{enumerate}
\end{shadedbox}
\end{column}
\begin{column}{0.5\textwidth}
\begin{shadedbox}[title=\textbf{What the Agent Can Do}]
\begin{itemize}\small
  \item Write and execute SQL queries
  \item Compute derived metrics (growth rates, ratios)
  \item Generate charts and export to Excel or PowerPoint
\end{itemize}
\end{shadedbox}
\end{column}
\end{columns}
\end{frame}

\begin{frame}{Replacing Dashboards: FP\&A and Treasury}
\begin{columns}[T]
\begin{column}{0.5\textwidth}
\begin{shadedbox}[title=\textbf{FP\&A}]
\begin{itemize}\small
  \item \textit{``Budget vs.\ actual for Q3, decompose the variance into volume, price, and cost drivers''}
  \item \textit{``Rolling 12-month revenue forecast with confidence bands''}
\end{itemize}
\end{shadedbox}
\end{column}
\begin{column}{0.5\textwidth}
\begin{shadedbox}[title=\textbf{Treasury \& Risk}]
\begin{itemize}\small
  \item \textit{``Cash position over the next 90 days using AR/AP forecasts''}
  \item \textit{``VaR by desk for the last 30 days, flag breaches''}
\end{itemize}
\end{shadedbox}
\end{column}
\end{columns}
\vspace{0.3cm}
Each of these would take a BI team \textbf{days to build}. With a database agent, they take \alert{seconds}.
\end{frame}

\begin{frame}{Replacing Dashboards: Portfolio and Executive}
\begin{columns}[T]
\begin{column}{0.5\textwidth}
\begin{shadedbox}[title=\textbf{Portfolio Management}]
\begin{itemize}\small
  \item \textit{``Sector allocation vs.\ benchmark with active weights''}
  \item \textit{``Performance attribution --- allocation vs.\ selection''}
\end{itemize}
\end{shadedbox}
\end{column}
\begin{column}{0.5\textwidth}
\begin{shadedbox}[title=\textbf{Executive Reporting}]
\begin{itemize}\small
  \item \textit{``One-page executive summary with KPIs and trends''}
  \item \textit{``Board deck from this quarter's financials --- 5 slides max''}
\end{itemize}
\end{shadedbox}
\end{column}
\end{columns}
\vspace{0.3cm}
The AI doesn't replace the analyst's judgment --- it replaces the \alert{mechanical work} of pulling data and building charts.
\end{frame}

\begin{frame}{The Weekly Ops Review --- Before and After}
\begin{columns}[T]
\begin{column}{0.5\textwidth}
\begin{shadedbox}[title=\textbf{Before: The Dashboard Era}]
\begin{enumerate}\small
  \item Data team pulls exports and builds slides (6+ hrs)
  \item Manager reviews and revises (3 hrs)
  \item VP asks a question --- ``We'll get back to you next week''
\end{enumerate}
\vspace{0.2cm}
\centering\small\textbf{Total: 9+ hours per week}
\end{shadedbox}
\end{column}
\begin{column}{0.5\textwidth}
\begin{shadedbox}[title=\textbf{After: Natural Language AI}]
\begin{enumerate}\small
  \item AI agent generates ops review (3 min)
  \item Manager reviews and iterates in chat
  \item VP asks a question --- \alert{AI answers in 10 seconds}
\end{enumerate}
\vspace{0.2cm}
\centering\small\textbf{Total: 15 minutes + live Q\&A}
\end{shadedbox}
\end{column}
\end{columns}
\end{frame}

% ========================================
% SECTION 6: THE REPORTING PIPELINE
% ========================================
\section{The Reporting Pipeline}

\begin{frame}{Five Steps from Query to Report}
\begin{center}
\begin{tikzpicture}[
  node distance=0.8cm,
  every node/.style={font=\small},
  sbox/.style={draw=titlegray, rounded corners, minimum width=2cm, minimum height=0.8cm, fill=excelinput, text=titlegray, font=\scriptsize\bfseries, align=center},
  sarrow/.style={-{Stealth[length=3mm]}, thick, draw=accentblue}
]
\node[sbox] (db) {Database\\Query};
\node[sbox, right=of db] (transform) {Transform\\\& Compute};
\node[sbox, right=of transform] (chart) {Charts \&\\Tables};
\node[sbox, right=of chart] (narr) {AI\\Narrative};
\node[sbox, right=of narr] (pptx) {PowerPoint\\Output};

\draw[sarrow] (db) -- (transform);
\draw[sarrow] (transform) -- (chart);
\draw[sarrow] (chart) -- (narr);
\draw[sarrow] (narr) -- (pptx);
\end{tikzpicture}
\end{center}
\vspace{0.5cm}
From raw data to polished deck --- \alert{no copy-paste, no formatting, no manual writing.}
\end{frame}

\begin{frame}{Building It with Your Coding Agent}
You do not need to write this by hand.  Give your coding agent a single prompt:
\vspace{0.3cm}
\begin{shadedbox}
\small\textit{``Build an app that connects to a database.  When the user selects a report and clicks Generate, run the SQL query, create a chart, write a three-sentence executive summary, and assemble a three-slide PowerPoint (title, chart, findings).  Show the chart and summary on screen with a download button for the PowerPoint file.''}
\end{shadedbox}
\vspace{0.3cm}
Your job is not to write the code.  Your job is to \alert{test the result, refine the prompt, and iterate} until the output meets your standards.
\end{frame}

\begin{frame}{Building the Pipeline with Streamlit}
\textbf{Streamlit} turns a Python script into a web app in minutes.  Combined with an AI pipeline, it becomes a \alert{self-service reporting tool}.
\vspace{0.3cm}
\begin{columns}[T]
\begin{column}{0.5\textwidth}
\begin{shadedbox}[title=\textbf{User Interface}]
\begin{itemize}\small
  \item Dropdown: select time period, department, or metric
  \item Button: ``Generate Report''
  \item Download: auto-generated \texttt{.pptx} file
\end{itemize}
\end{shadedbox}
\end{column}
\begin{column}{0.5\textwidth}
\begin{shadedbox}[title=\textbf{Behind the Scenes}]
\begin{itemize}\small
  \item Query database for selected parameters
  \item Compute metrics and build charts
  \item Send chart data to LLM for narrative
  \item Assemble PowerPoint with \texttt{python-pptx}
\end{itemize}
\end{shadedbox}
\end{column}
\end{columns}
\vspace{0.3cm}
The user clicks one button and gets a polished deck.  \alert{No analyst needed.}
\end{frame}

\begin{frame}{Example: Monthly Ops Review Pipeline}
\begin{columns}[T]
\begin{column}{0.5\textwidth}
\textbf{What the Pipeline Produces}
\begin{enumerate}\small
  \item Title slide, KPI summary table, and trend charts
  \item Variance waterfall (budget vs.\ actual)
  \item AI-generated executive summary
\end{enumerate}
\end{column}
\begin{column}{0.5\textwidth}
\textbf{Time Comparison}
\begin{itemize}\small
  \item \textbf{Manual}: 6--8 hours per month (data pull, Excel, copy-paste into PowerPoint, write narrative)
  \item \textbf{Pipeline}: \alert{30 seconds} per run
  \item \textbf{ROI}: First month pays for the build time
\end{itemize}
\end{column}
\end{columns}
\end{frame}

\begin{frame}[fragile]{Automated PowerPoint with python-pptx}
\texttt{python-pptx} is a Python library for creating and editing PowerPoint files programmatically.
\vspace{0.3cm}
\begin{columns}[T]
\begin{column}{0.5\textwidth}
\begin{shadedbox}[title=\textbf{What It Can Do}]
\begin{itemize}\small
  \item Create slides from templates
  \item Insert charts, tables, and images
  \item Apply corporate formatting (fonts, colors, logos)
  \item Populate placeholders with live data
\end{itemize}
\end{shadedbox}
\end{column}
\begin{column}{0.5\textwidth}
\begin{shadedbox}[title=\textbf{The AI Advantage}]
\begin{itemize}\small
  \item AI writes the \texttt{python-pptx} code for you
  \item Tell Claude: ``Create a 5-slide deck from this data with a waterfall chart on slide 3''
  \item AI handles layout, formatting, and data binding
\end{itemize}
\end{shadedbox}
\end{column}
\end{columns}
\vspace{0.3cm}
\alert{You describe the deck; AI builds the automation.}  The pipeline runs unattended after that.
\end{frame}

% ========================================
% SECTION 7: SANDBOXED EXECUTION
% ========================================
\section{Sandboxed Execution}

\begin{frame}{Development vs.\ Production}
\begin{columns}[T]
\begin{column}{0.5\textwidth}
\begin{shadedbox}[title=\textbf{Development (Your Laptop)}]
\begin{itemize}\small
  \item Agent runs in your environment
  \item Full access to your files
  \item Fine for prototyping
\end{itemize}
\end{shadedbox}
\end{column}
\begin{column}{0.5\textwidth}
\begin{shadedbox}[title=\textbf{Production}]
\begin{itemize}\small
  \item Agent runs in a \alert{container} (Docker)
  \item Isolated, disposable environment
  \item Read-only database access
\end{itemize}
\end{shadedbox}
\end{column}
\end{columns}
\vspace{0.3cm}
A container is a disposable, isolated computing environment.  A bug in AI-generated code cannot affect your other files, your database, or your network.
\end{frame}

\begin{frame}{The Sandbox Pattern}
\begin{center}
\begin{tikzpicture}[
  node distance=0.8cm and 1.5cm,
  every node/.style={font=\small},
  sbox/.style={draw=titlegray, rounded corners, minimum width=2.2cm, minimum height=0.9cm, fill=excelinput, text=titlegray, font=\scriptsize\bfseries, align=center},
  sarrow/.style={-{Stealth[length=2.5mm]}, thick, draw=accentblue}
]
\node[sbox] (user) {User};
\node[sbox, right=1.2cm of user] (ui) {Web UI};
\node[sbox, right=1.2cm of ui] (agent) {Agent};
\node[sbox, right=1.2cm of agent, fill=alertorange!15, draw=alertorange] (code) {Sandboxed\\Code};
\node[sbox, below=0.8cm of code] (db) {Database\\{\tiny read-only}};
\node[sbox, below=0.8cm of agent] (report) {Report};

\draw[sarrow] (user) -- (ui);
\draw[sarrow] (ui) -- (agent);
\draw[sarrow] (agent) -- (code);
\draw[sarrow] (code) -- (db);
\draw[sarrow] (db) -- (code);
\draw[sarrow] (code) -- (report);
\draw[sarrow] (report) -- (ui);
\end{tikzpicture}
\end{center}
\vspace{0.2cm}
\begin{itemize}
  \item User requests a report through the web interface
  \item Agent plans and orchestrates; code runs in a sandboxed container
  \item Container queries the database through a \alert{read-only} connection
  \item Finished report flows back to the user
\end{itemize}
\end{frame}

\begin{frame}{From Prototype to Production}
The reporting app on your laptop is a prototype.  Moving it to production adds \textbf{infrastructure}, not intelligence.
\vspace{0.3cm}
\begin{itemize}
  \item \textbf{Authentication}: SSO so only authorized users can generate reports
  \item \textbf{Scheduling}: weekly reports generated automatically
  \item \textbf{Logging}: every query and LLM call recorded for compliance
  \item \textbf{Read-only credentials}: the agent cannot modify data
\end{itemize}
\vspace{0.3cm}
The pipeline you prototyped is the same artifact that powers the production tool.  IT wraps it in infrastructure; \alert{you provide the domain knowledge and the prompt}.
\end{frame}

% ========================================
% SECTION 8: VARIANCE ANALYSIS
% ========================================
\section{Variance Analysis}

\begin{frame}{What is Variance Analysis?}
\textbf{Variance analysis} compares \alert{budgeted} figures to \alert{actual} results and decomposes the differences into actionable drivers. It is the core analytical task in FP\&A.
\vspace{0.3cm}
\begin{columns}[T]
\begin{column}{0.5\textwidth}
\textbf{Why It Matters}
\begin{itemize}\small
  \item Every public company does it quarterly
  \item Boards and investors ask ``why did we miss?''
  \item Drives re-forecasting and capital allocation
\end{itemize}
\end{column}
\begin{column}{0.5\textwidth}
\textbf{Key Decompositions}
\begin{itemize}\small
  \item \textbf{Revenue}: volume vs.\ price vs.\ mix
  \item \textbf{COGS}: volume vs.\ unit cost
  \item \textbf{SG\&A}: headcount vs.\ rate vs.\ discretionary
\end{itemize}
\end{column}
\end{columns}
\vspace{0.3cm}
Revenue variance = \alert{Volume effect} + \alert{Price effect} + \alert{Mix effect}
\end{frame}

\begin{frame}{Variance Analysis: Example Data}
\textbf{Scenario:} You are an FP\&A analyst. Q1 actuals just closed. The CEO wants to know why operating income missed budget by \$300K.
\vspace{0.3cm}
\begin{itemize}
  \item \textbf{Revenue:} Budget 100K units at \$50 = \$5M; Actual 95K units at \$51 = \$4.845M \quad \alert{(\$155K) miss}
  \item \textbf{COGS:} Budget \$30/unit = \$3M; Actual \$32/unit = \$3.04M \quad \alert{(\$40K) miss}
  \item \textbf{SG\&A:} Budget \$1.45M; Actual \$1.555M \quad \alert{(\$105K) miss}
  \item \textbf{Operating Income:} Budget \$550K; Actual \$250K \quad \alert{(\$300K) total miss}
\end{itemize}
\end{frame}

\begin{frame}[shrink=5]{Variance Analysis with an Agent}
With a database agent connected to your financials, one prompt replaces 2--4 hours of FP\&A work:
\vspace{0.2cm}

\textbf{Prompt}\\[0.3em]
\small ``Here is our Q1 budget vs.\ actuals spreadsheet. Decompose the \$300K operating income miss into volume, price, cost, and discretionary spending drivers.  Produce a waterfall chart and a summary memo for the CFO.''
\normalsize
\vspace{0.2cm}

\textbf{What the Agent Does}
\begin{enumerate}\small
  \item \textbf{Plans}: identifies the variance decomposition formulas needed
  \item \textbf{Decomposes} revenue into volume + price; COGS into volume + rate
  \item \textbf{Generates} waterfall chart and writes CFO-ready memo in Word format
\end{enumerate}
\end{frame}

% ========================================
% FINANCE APPLICATION: M&A DUE DILIGENCE
% ========================================
\section{Finance Application: M\&A Due Diligence}

\begin{frame}{M\&A Due Diligence with an Agent}
Give an agent a goal: ``Evaluate this acquisition target.'' The agent:
\vspace{0.3cm}
\begin{enumerate}
  \item Ingests data from multiple formats (Excel, PDF, Word, CSV)
  \item Applies evaluation criteria and computes risk metrics
  \item Produces a summary memo with flagged risks
\end{enumerate}
\vspace{0.3cm}
\alert{This is a single prompt to Claude Code.} The agent plans and executes all three steps autonomously, reading multiple file formats and combining the results.
\end{frame}

% ========================================
% THE ORCHESTRATION LAYER
% ========================================
\section{Orchestration}

\begin{frame}{The Orchestration Layer}
The agent's control logic can route different tasks to different models and prompts.
\vspace{0.3cm}
\begin{columns}[T]
\begin{column}{0.5\textwidth}
\begin{shadedbox}[title=\textbf{Different Prompts}]
\begin{itemize}\small
  \item SQL generation $\rightarrow$ database schema prompt
  \item Python analysis $\rightarrow$ data science prompt
  \item Each task gets specialized instructions
\end{itemize}
\end{shadedbox}
\end{column}
\begin{column}{0.5\textwidth}
\begin{shadedbox}[title=\textbf{Different Models}]
\begin{itemize}\small
  \item Simple classification $\rightarrow$ fast, cheap model
  \item Complex reasoning $\rightarrow$ powerful model
  \item Cost and speed optimization
\end{itemize}
\end{shadedbox}
\end{column}
\end{columns}
\vspace{0.3cm}
Sub-agents: dispatch specialized workers for parallel tasks. This is what Cowork does with its parallel VMs.
\end{frame}

% ========================================
% EXERCISES
% ========================================
\section{Exercises}

\begin{frame}{Exercise 1: Multi-Step Agent Workflow}
Give Claude Code a single compound instruction that requires at least three steps.
\vspace{0.3cm}
\textbf{Example:} \textit{``Fetch Apple's quarterly revenue from the Rice Data Portal.  Create a bar chart of the last 8 quarters.  Write a one-paragraph executive summary.  Save the chart and summary to a reports folder.''}
\vspace{0.3cm}
\begin{itemize}
  \item Observe the agent loop --- how many distinct tool calls does the agent make?
  \item Does it check its own work?
  \item If the output is not right, refine your prompt and try again
\end{itemize}
\end{frame}

\begin{frame}{Exercise 2: Workflow Decomposition}
Consider the task: \textit{``Prepare the quarterly business review for the CEO.''}
\vspace{0.3cm}
\begin{enumerate}
  \item List the 5--7 steps a human analyst would take to complete this task
  \item For each step, identify what \textbf{tool} the agent would use and what \textbf{data} it would need
  \item Mark which steps need \alert{human approval} before the agent continues
  \item Write a single prompt that describes the full workflow for an agent
\end{enumerate}
\end{frame}

\begin{frame}{Exercise 3: M\&A Due Diligence}
\begin{enumerate}
  \item Download the due diligence data pack (Excel + PDF + Word + CSV)
  \item Ask Claude Code to evaluate the acquisition target end-to-end
  \item Submit: the summary report + screenshots of intermediate steps
\end{enumerate}
\vspace{0.3cm}
Watch how Claude plans its approach, reads each file, and combines the results into a coherent analysis.
\end{frame}

\begin{frame}{Exercise 4: Streamlit App with the Anthropic API}
\begin{enumerate}
  \item Get an API key from \texttt{console.anthropic.com}
  \item Ask Claude Code to build a Streamlit app that:
  \begin{itemize}\small
    \item Takes a company ticker from the user
    \item Sends a prompt to the Anthropic API asking for an investment summary
    \item Displays the AI-generated summary on screen
  \end{itemize}
  \item Extend: add a system prompt with specific analysis criteria (valuation, growth, risks)
\end{enumerate}
\end{frame}

\end{document}
